\clearpage
\thispagestyle{empty}
\begin{center}
    \bfseries\fontsize{14}{16}\selectfont RESUMEN
\end{center}
\vspace{16pt}

\noindent
El \textit{Capacitated Vehicle Routing Problem} (CVRP) es un problema clásico de
optimización combinatoria en el que se debe determinar un conjunto de rutas de costo
mínimo para abastecer clientes con demandas conocidas mediante una flota homogénea de
vehículos con capacidad limitada. En este informe se aborda el CVRP
exclusivamente desde la perspectiva de los métodos exactos de programación matemática.
Se revisa el panorama de algoritmos \textit{Branch and X} para el problema, abarcando
\textit{Branch and Bound}, \textit{Branch and Cut}, \textit{Branch and Price} y
\textit{Branch, Cut and Price}, y se justifica la adopción del esquema clásico de
Laporte y Nobert (1983) como base metodológica del trabajo. El documento presenta la
definición formal del CVRP, su formulación matemática como programa lineal entero sobre
un grafo no dirigido, y procedimiento de resolución basado en la
generación iterativa de desigualdades de capacidad sobre una relajación inicial con
restricciones de grado. La implementación computacional, desarrollada en Python mediante PuLP y el
solver COIN-OR CBC, fue evaluada sobre cinco instancias estándar de CVRPLIB
(Sets E y A). El método certificó el óptimo en cuatro de ellas (gap 0\,\%) y,
en la quinta, alcanzó una solución factible a 0{,}91\,\% del óptimo dentro del
presupuesto de tiempo, evidenciando empíricamente el límite de escalabilidad
del esquema clásico.

\vspace{10pt}
\noindent\textbf{Palabras clave:} CVRP; optimización combinatoria; programación lineal
entera; métodos exactos; \textit{Branch and Bound}; CVRPLIB.